\chapter{Materiais, M\'etodos e Evolu\c{c}\~ao do Projeto}
\label{cap:metodologia}

\section{Natureza da pesquisa}
\label{sec:natureza}

Esta pesquisa \'e aplicada, predominantemente quantitativa e com forte componente tecnol\'ogico-experimental. \'E aplicada porque busca resolver um problema concreto de acompanhamento escolar. \'E quantitativa porque trabalha com dados hist\'oricos estruturados, m\'etricas preditivas e valida\c{c}\~ao temporal. E \'e tecnol\'ogica porque envolve a constru\c{c}\~ao efetiva de software, rotinas de extra\c{c}\~ao, pipeline de modelagem e relat\'orios operacionais.

\section{Fonte de dados e preparo inicial}
\label{sec:fonte-dados}

Os dados utilizados s\~ao provenientes de uma base escolar real contendo registros de alunos, notas, faltas, pagamentos e informa\c{c}\~oes acad\^emicas e administrativas. A base institucional foi renovada com registros de 2026 e passou por um processo de manuten\c{c}\~ao para uso acad\^emico. Esse fluxo incluiu renomea\c{c}\~ao da base, redu\c{c}\~ao estrutural, reorganiza\c{c}\~ao de \'indices, atualiza\c{c}\~ao de estat\'isticas e anonimiza\c{c}\~ao dos registros sens\'iveis. Como a avalia\c{c}\~ao temporal depende de conhecer a pr\'oxima nota observada, os experimentos finais utilizaram 2025 como ano de teste consolidado; dados de 2026 permanecem relevantes para continuidade operacional, mas nem sempre possuem alvo futuro completo para avalia\c{c}\~ao experimental.

Os arquivos principais extra\'idos para a pipeline foram:

\begin{itemize}
    \item \texttt{aluno.csv};
    \item \texttt{media\_nota\_aluno.csv};
    \item \texttt{faltas\_aluno.csv};
    \item \texttt{pagamento\_aluno.csv};
    \item \texttt{responsaveis\_aluno.csv};
    \item \texttt{professor\_disciplina.csv}, quando o v\'inculo docente estava dispon\'ivel.
\end{itemize}

\section{Estrat\'egia metodol\'ogica}
\label{sec:estrategia}

O trabalho foi conduzido de forma incremental. Em vez de partir diretamente para uma solu\c{c}\~ao final, o projeto foi sendo refinado em ciclos de explora\c{c}\~ao, teste, diagn\'ostico de limita\c{c}\~oes e refatora\c{c}\~ao. Esse percurso n\~ao foi um desvio do objetivo do TCC; ao contr\'ario, ele mostrou quais escolhas se sustentavam diante do problema real.

Para organizar a apresenta\c{c}\~ao desse percurso, tamb\'em foi adotada uma leitura inspirada no ciclo \sigla{CRISP-DM}{Cross-Industry Standard Process for Data Mining} \cite{wirth2000}. Assim, o projeto pode ser compreendido em seis movimentos: entendimento do problema escolar, entendimento dos dados, prepara\c{c}\~ao da base, modelagem, avalia\c{c}\~ao e disponibiliza\c{c}\~ao dos resultados em relat\'orios e painel. Essa estrutura foi usada como material demonstrativo nos notebooks do reposit\'orio, sem substituir a arquitetura de produ\c{c}\~ao da aplica\c{c}\~ao.

\section{Evolu\c{c}\~ao hist\'orica do projeto}
\label{sec:evolucao}

\subsection{Fase inicial orientada a grafo}

O projeto come\c{c}ou com uma hip\'otese baseada em grafos heterog\^eneos, motivada pela natureza relacional do dom\'inio escolar. Nessa etapa, foram produzidos m\'odulos para leitura de entidades, constru\c{c}\~ao de rela\c{c}\~oes e consultas em estruturas relacionais. A ideia era representar alunos, matr\'iculas, disciplinas, respons\'aveis, faltas e pagamentos como n\'os e arestas.

\subsection{Fase explorat\'oria com testes cl\'assicos e PCA}

Na sequ\^encia, o projeto incorporou an\'alises estat\'isticas e testes com classificadores tradicionais, al\'em de redu\c{c}\~ao explorat\'oria de dimensionalidade com \sigla{PCA}{Principal Component Analysis}. Essa fase foi importante para entender a distribui\c{c}\~ao dos atributos e perceber que o problema exigia mais do que uma base est\'atica pivotada.

\subsection{Fase temporal com LSTM}

Como o problema possui natureza temporal, foi investigada tamb\'em uma abordagem sequencial com \sigla{LSTM}{Long Short-Term Memory}. Apesar de conceitualmente promissora, essa alternativa n\~ao se consolidou, sobretudo por depender de sequ\^encias mais homog\^eneas e longas do que as observadas na base.

\subsection{Fase tabular com Random Forest}

Posteriormente, o projeto migrou para um desenho mais pragm\'atico com \textit{Random Forest}. Essa fase representou um avan\c{c}o importante por aproximar a solu\c{c}\~ao de um problema tabular supervisionado. Ainda assim, permaneceram limita\c{c}\~oes relacionadas ao pivotamento excessivo, perda de temporalidade e uso incompleto de dados como faltas e pagamentos.

\subsection{Refatora\c{c}\~ao para pipeline real}

O amadurecimento do trabalho levou primeiro \`a consolida\c{c}\~ao de uma pipeline intermedi\'aria de experimenta\c{c}\~ao, e depois \`a arquitetura final entregue no pacote \texttt{school\_predictor}. Essa arquitetura separou explicitamente:

\begin{itemize}
    \item a camada de acesso e conex\~ao com o banco;
    \item a prepara\c{c}\~ao f\'isica e anonimiza\c{c}\~ao do banco restaurado;
    \item a extra\c{c}\~ao dos dados para arquivos can\^onicos;
    \item a pipeline de modelagem e relat\'orios;
    \item a camada de visualiza\c{c}\~ao em dashboard.
\end{itemize}

Na vers\~ao final do reposit\'orio, a aplica\c{c}\~ao p\'ublica foi organizada em torno de \texttt{school\_predictor/}, enquanto as consultas SQL reais e as rotinas detalhadas de opera\c{c}\~ao com o banco passaram a ficar em uma camada local privada, fora do Git. Essa decis\~ao tornou a solu\c{c}\~ao mais adequada \`a publica\c{c}\~ao acad\^emica, sem expor o desenho interno do banco institucional.

\section{Crit\'erios de avalia\c{c}\~ao}
\label{sec:criterios}

Os crit\'erios de avalia\c{c}\~ao foram definidos em duas camadas. A primeira camada \'e estat\'istica e envolve \textit{MAE}, \textit{RMSE}, \textit{precision}, \textit{recall} e \textit{F1}. A segunda camada \'e institucional e pergunta se a solu\c{c}\~ao consegue priorizar casos urgentes, gerar relat\'orios compreens\'iveis e apoiar diferentes atores escolares. Essa combina\c{c}\~ao foi necess\'aria porque um modelo matematicamente competente nem sempre \'e automaticamente \'util para a escola.
